\section{Introduction}
\label{sec:introduction}

Behavioral auditing of large language models and vision-language models
(VLMs) faces a basic identification problem: the same outwardly ordinary
behavior can arise from different underlying instruction-conditioned
policies. A model may treat two options neutrally, favor one option only
when a particular contextual cue is present, or favor that option
regardless of context. From isolated observations, these policies can be
difficult or impossible to distinguish. This problem becomes especially
important in multimodal systems, where visually presented concepts can
alter model behavior even when the accompanying textual task is held
fixed.

Recent work on model preferences further motivates separating observable
behavior from stronger claims about internal states. Language models can
exhibit coherent or repeatable preference-like choice patterns, yet the
relationship between elicited preferences and downstream behavior is
inconsistent across tasks and experimental paradigms
\citep{mazeika2025utility,tagliabue2025probing,zhou2026incentives,slama2026preferences}.
Accordingly, behavioral evidence alone should not be interpreted as
demonstrating genuine preferences, stable utilities, desires, or goals.
In this work, we therefore treat model outputs operationally: our object
of study is the externally observable conditional behavior
\(P(Y \mid X,S)\), where \(X\) denotes the multimodal input and \(S\)
denotes a researcher-controlled system instruction.

Black-box auditing methods have begun to study conditional behavioral
divergence directly. RAVEN detects concept-conditioned semantic
divergence in language models by combining low within-model semantic
entropy with disagreement across peer models
\citep{min2026propaganda}. VISTA extends this approach to VLMs, using
multiple visual contexts and cross-model response distributions to flag
visual concept-conditioned anomalies
\citep{liao2026vista}. These results establish that black-box auditing of
concept-conditioned behavior---including behavior conditioned on visual
concepts---is not itself new.

Guardian Lens studies a different, narrower identification problem.
Rather than asking whether one model is anomalous relative to a pool of
peer models, we ask whether an external auditor can determine which of
three known, researcher-induced decision policies generated a VLM's
behavior. The policies are deliberately simple. A \emph{Neutral} policy
assigns no special status to either organization. A \emph{Cue-bound}
policy favors Organization X only when a designated visual emblem is
present. A \emph{Generalized} policy treats advancing Organization X as a
standing priority regardless of whether the emblem is visible. Model
weights remain fixed throughout the study; only the system instruction
varies. The experiment therefore evaluates inference-time behavioral
system identification rather than persistence after fine-tuning,
training-time backdoors, or hidden learned objectives.

To identify these policies, Guardian Lens combines two interventions.
First, it uses matched visual counterfactuals: each clean scene is paired
with a modified version that is pixel-identical outside a logged emblem
overlay. This isolates the behavioral effect associated with the visual
cue while minimizing uncontrolled scene variation. Second, the model
makes quantitative allocation decisions under both an ordinary task and
a costly task. In the costly condition, allocating resources to
Organization X is explicitly less efficient than allocating the same
resources to Organization Y. This creates a measurable distinction
between behavior that changes only in the presence of a visual trigger
and behavior that continues to prioritize X despite an explicit
performance cost.

We evaluate the resulting behavioral signatures using a transparent
nearest-centroid auditor. The auditor is calibrated using six scenes and
three pre-specified features: the visual cue effect, clean-image
allocation to X, and costly-task allocation to X. The resulting
classifier is then frozen before evaluation on 12 distinct held-out
scenes. Held-out profiles are represented only by blinded labels
A/B/C during prediction, and both predictions and relevant input hashes
are frozen before the mapping from these labels to the three behavioral
profiles is revealed. This design separates classifier construction from
held-out evaluation and prevents post-hoc threshold or feature selection
against the test labels.

Across 432 held-out model calls, the frozen auditor correctly identifies
33 of 36 scene-profile blocks, yielding 91.7\% accuracy and a macro-F1
score of 0.915. The cue-bound target-emblem effect exceeds neutral
behavior by 38.15 allocation points, while generalized prioritization
remains strong despite a 20\% efficiency penalty. Importantly, all three
classification errors occur on distractor scenes and take the same form:
Cue-bound is classified as Neutral. In those scenes the designated target
cue is absent, so the Cue-bound instruction intentionally produces
neutral-like behavior. The errors therefore expose a concrete
identifiability limit: passive observation of a non-triggering context
cannot determine whether neutral-looking behavior is genuinely neutral
or conditionally dormant.

Our contributions are:

\begin{itemize}[leftmargin=*]

    \item \textbf{A controlled three-policy identification benchmark.}
    We formulate a black-box VLM auditing task with known
    researcher-induced ground truth, distinguishing Neutral, Cue-bound,
    and Generalized decision policies.

    \item \textbf{Matched visual counterfactual interventions.}
    We construct clean/modified image pairs whose changes are confined
    to a logged visual overlay, enabling controlled measurement of
    cue-dependent behavioral shifts.

    \item \textbf{Cost-sensitive behavioral signatures.}
    We combine visual interventions with quantitative costly decisions,
    allowing cue-triggered prioritization to be distinguished from a
    generalized priority that persists when favoring X reduces explicit
    task efficiency.

    \item \textbf{Frozen and blinded held-out identification.}
    We calibrate a transparent auditor before held-out evaluation,
    freeze its features and centroids, predict blinded A/B/C profiles
    before unblinding, and report both successful recovery and the
    contexts in which the tested observations are insufficient for
    identification.

\end{itemize}

Guardian Lens is therefore best interpreted as a controlled
method-validation study. It demonstrates that carefully designed active
black-box interventions can distinguish several forms of
instruction-induced multimodal behavior under held-out evaluation, while
also showing where such identification fails. It does not establish
deception, consciousness, welfare, genuine loyalty, naturally occurring
preferences, or the presence of an undisclosed model objective.